\section{海水耦合瞬态热场与守恒降阶演化}\label{sec:thermal}
\subsection{固体与海水的热方程}
对封装、导体、磁芯等固体，在自身材料坐标内
\begin{equation}
\rho_i c_i(T)\,\mathrm D_t T_i
=\grad\cdot(k_i(T)\grad T_i)+q_i^{\rm EM}.
\label{eq:heat-solid}
\end{equation}
对于海水
\begin{equation}
\rho_sc_s(T_s)[\partial_tT_s+\bm v_s\cdot\grad T_s]
=\grad\cdot(k_s(T_s)\grad T_s)+q_s^{\rm EM}.
\label{eq:heat-sea}
\end{equation}
界面给出接触温差/接触热阻及热通量平衡；
海水远场具有背景温度/热输运条件。
第一阶段可采用静止海水或指定流速场，
但不能宣称不建流场就能正确预测任意湍流与自然对流。
运动组件中的温度采用材料导数，避免将刚体运输重复计入热源。

\subsection{几何条件化连续温度场}
对第 \(i\) 个组件使用随体坐标
\(\bm\xi_i=R_i^{\mathsf T}(\bm x-\bm c_i)\)，
研究连续低维表达
\begin{equation}
T_i(\bm x,t)=T_{\rm ref}(\bm x,t)
+\sum_{\ell=1}^{r_i}
a_{i,\ell}(t)\Phi_{i,\ell}(\bm\xi_i;\bm g_i).
\label{eq:thermal-basis}
\end{equation}
\(\Phi\) 可来自解析 Green/传热本征基与有限维学习修正，
需对边界接口和热容量保持物理一致。
\textbf{不可默认固定 FEM 网格 POD 基支持全部 BFZI 形状与盒体尺寸}。
海水的热尾迹应使用独立的连续核/低阶状态表达，
否则可能错误丢失流向下游的热能。

\subsection{守恒低阶热动力系统}
对热量方程进行几何条件化投影，得到候选
\begin{equation}
\boxed{
\bm M(\bm g,\bm a)\dot{\bm a}
+\bm K(\bm g,\bm a,\bm q,\bm v_s)\bm a
=\bm B(\bm g)\bm p_{\rm EM}
+\bm f_\theta(\bm a,\bm g,\bm q,\bm v_s).
}
\label{eq:thermal-rom}
\end{equation}
这里 \(\bm M\succ0\)，
\(\bm K\) 对应传导/换热耗散，
学习闭合项 \(\bm f_\theta\) 只能表示未建模热通量
或低阶截断，不得任意制造/消灭能量。
在线时间通过对式 \eqref{eq:thermal-rom} 的积分推进，
不是将时间作为单个黑箱网络输入后截断训练窗。
这给出长时间求解的机制，\textbf{并不自动证明长时间精度}；
仍需在未知功率历程、位姿和热稳态下独立测试。
